\documentclass{strategynet-report}

\hypersetup{
  pdftitle={Slice + Hypercube},
  pdfauthor={research@strategynet.ai}
}

\lstdefinestyle{slice}{
  style=sncode
}

\newcommand{\SLICE}{\textsf{Slice}}
\newcommand{\R}{\mathbb{R}}
\newcommand{\E}{\mathbb{E}}

\tikzset{
  slicebox/.style={
    draw=SNInk,
    fill=SNInk!5,
    rounded corners=2pt,
    align=center,
    minimum height=9mm,
    text width=29mm,
    font=\sffamily\fontsize{8}{9.5}\selectfont
  },
  slicehot/.style={
    slicebox,
    draw=SNSignalDark,
    fill=SNSignal!7
  },
  slicemeta/.style={
    slicebox,
    draw=SNMuted,
    fill=SNMuted!7
  },
  sliceaction/.style={
    slicebox,
    draw=SNViolet,
    fill=SNViolet!7
  },
  sliceflow/.style={
    -{Latex[length=2mm]},
    draw=SNInk,
    line width=0.75pt
  },
  slicefuture/.style={
    -{Latex[length=2mm]},
    draw=SNSignalDark,
    dashed,
    line width=0.75pt
  }
}

\title{\SLICE{} + Hypercube}
\SNSubtitle{Current analytical state without rebuilding it for every reader}
\author{\href{mailto:research@strategynet.ai}{research@strategynet.ai}}
\SNDocumentType{Foundations and architecture}
\SNVersion{Version 0.2}
\SNDate{July 2026}
\SNStatus{Open-source working paper}
\SNShortTitle{Slice + Hypercube}

\begin{document}
\maketitle

\begin{abstract}
A live basket monitor may recalculate thousands of values whenever prices
change.  The arithmetic is usually cheap.  Repeatedly finding, joining,
serializing, and copying the same current inputs is not.  When every reader
rebuilds that state for itself, a simple vector calculation inherits database
latency, extra allocation, and inconsistent observation boundaries.

\SLICE{} gives local processes a typed, memory-mapped view of current,
entity-aligned values.  Hypercube evaluates one coherent input generation
through a declared node graph and produces cross-sectional outputs that can be
published back as aligned slices.  The public repository includes the format,
readers and writers, calculation engine, deterministic synthetic sources, and
live browser and terminal demonstrations.  This paper explains that
implemented path and its limits.  Separate slices are not yet one atomic
multi-slice generation; vendor feeds, durable persistence, arbitrary function
cells, replication, GPU lowering, and trading actions remain outside the
current release.  Although the examples use markets, the API also works for
any stable set of entities.
\end{abstract}

\SNtableofcontents

\part{Motivation}

\section{From events to live analytical state}

Markets arrive as events, but many decisions are functions of state.  A quote
event changes the best bid and offer for one instrument; the next consumer
usually wants the latest midpoint, spread, return, risk exposure, or
cross-sectional rank.  If every consumer independently rebuilds that state
from an event stream, the system duplicates joins, symbol resolution, window
maintenance, and normalization.  If every calculation first queries a durable
database, the database becomes part of the latency and availability path even
when the working set is only a few megabytes of current vectors.

The distinction is temporal rather than ideological:
\begin{itemize}
  \item an event log answers what happened and in what order;
  \item a durable analytical store answers historical and ad hoc questions;
  \item a live state plane answers what the current published values are;
  \item an execution graph answers how declared outputs follow from those
        values.
\end{itemize}
These are complementary responsibilities.  \SLICE{} is the live state plane.
The hypercube is the logical coordinate system and calculation graph.  A
database remains the recovery, research, and audit path; a transport remains
the mechanism for moving events or published generations between machines.

\begin{SNthesis}
A current analytical value should be cheap to read because its identity,
layout, type, generation, and provenance were established when it was
published---not rediscovered by every reader.
\end{SNthesis}

The architecture is motivated by repeated whole-population work.  Given
entity-aligned weights \(w\), returns \(r\), signals \(s\), and betas
\(\beta\), common calculations are
\[
  R_p=w^\top r,\qquad
  E_s=w^\top s,\qquad
  E_\beta=w^\top\beta,\qquad
  G=\lVert w\rVert_1.
\]
The arithmetic is elementary.  Most accidental complexity lies in obtaining
vectors that refer to the same entities, units, observation boundary, and
generation.  The system therefore treats alignment and publication semantics
as first-class data contracts.

\section{Design goals}

The present architecture has six goals:
\begin{enumerate}
  \item maintain latest source and derived state in fixed-layout,
        memory-mapped vectors;
  \item permit many local processes to perform point reads and vector scans
        without taking a shared reader mutex;
  \item reject linear algebra across incompatible entity layouts or value
        schemas;
  \item run the same node graph in live, replay, and batch modes by changing
        the supplied update rather than the calculation;
  \item publish node identity, observation timestamps, coverage, and local
        compute diagnostics beside values;
  \item leave room for typed, composable function cells without
        allowing calculation code to inherit trading authority.
\end{enumerate}

The first implementation intentionally narrows the platform contract.  It is
little-endian, uses fixed-size pointer-free payloads, and assumes one writer
with many readers per slice.  Those constraints make the shared ABI explicit
enough for Rust today and C++ or Python readers later.  A richer type system,
distributed replication, global multi-slice generations, and effectful
actions are evolution points, not implicit properties of the current files.

\section{Related work and scope}

Memory mapping is the operating-system foundation.  POSIX specifies a mapping
between a process address space and a file or shared-memory object and defines
shared mappings through which writes change the underlying object
\cite{posixmmap}.  \SLICE{} adds an application ABI, identity, layout, and
consistency protocol to that primitive.  It does not assume that mapping a
file alone makes concurrent access coherent.

Columnar memory systems establish the advantages of contiguous, typed buffers.
Apache Arrow defines a language-independent columnar representation designed
for sequential scans, constant-time random access, vectorization, and
relocatable zero-copy shared-memory access \cite{arrow}.  \SLICE{} shares
those physical motivations but has a smaller purpose: mutable current-state
vectors with stable entity slots and explicit writer epochs.  It is not an
Arrow implementation or a replacement for its interoperability ecosystem.

Array database work treats multidimensional arrays as first-class data rather
than encoding every analytical object as a relation.  SciDB, for example, was
designed for complex analytics over array-shaped scientific data
\cite{scidb}.  The hypercube developed here is closer to a resident execution
and publication model than a general array DBMS.  Dimensions describe aligned
live slices and calculation identities; durable historical queries remain
outside the hot path.

The numerical interface follows a long tradition of separating stable vector
and matrix operations from their machine-specific implementations.  The
original BLAS specification included dot products, vector operations, norms,
and related kernels \cite{blas}.  \SLICE{} does not define a new numerical
algebra.  It makes aligned live vectors safe to hand to such an algebra and
leaves room for optimized CPU, SIMD, or GPU kernels behind the same logical
operation.

For change propagation, adaptive functional programming and differential
dataflow show how explicit dependencies allow a system to update affected
results rather than recompute an entire program \cite{adaptive,differential}.
Those systems are broader and more formal than the current hypercube engine.
They nevertheless supply the right conceptual boundary for future function
cells: a declared dependency graph, logical time, deterministic propagation,
and controlled state are preferable to hidden callbacks attached to mutable
memory.

\part{Memory and Cube Model}

\section{Stable entity layout}

Let a layout be
\[
  L=(\lambda,a,n,U,\iota),
\]
where \(\lambda\) is a layout identity, \(a\) is a normalized namespace
(retained as \texttt{asset\_class} in version-1 JSON), \(n\) is allocated
capacity, \(U\) is the enabled entity set, and
\(\iota:U\rightarrow\{0,\ldots,n-1\}\) is an injective slot mapping.  A symbol
is a normalized lookup key; \texttt{instrument\_id} is the legacy field name
for the stable namespaced entity identity, and the slot is the dense numerical
address.

\[
  \text{entity key}\longrightarrow\text{stable entity id}
  \longrightarrow\text{slot id}\longrightarrow S[\text{slot id}].
\]

Slots do not move within a layout.  New entities consume unused slots;
removed entities are disabled or tombstoned.  Reordering or compaction
creates a new layout.  The serialized layout is hashed, and every dependent
slice carries that hash.  A dot product is rejected unless its inputs agree on
layout hash, schema, and active length.

This separation matters because a domain registry and a numerical layout
have different rates of change.  Venue, currency, multiplier, expiry, strike,
and other rich metadata can live in keyed sidecars or services.  The layout
keeps only the identity and stable dense index required by the shared vectors.

\section{Slice as a published vector}

For a value field \(f\) and a published generation \(g\), a primitive slice is
a partial vector
\[
  S_{f,g}:U\rightarrow V_f\cup\{\bot\},
\]
where \(V_f\) is a fixed physical type and \(\bot\) denotes missing or invalid
state.  Dense analytical slices usually use \(V_f=\R\); point-state slices use
fixed records such as quote, trade, or quote-at-trade.  Validity is explicit
rather than inferred from an arbitrary zero.

The version-one file begins with a 256-byte header:
\begin{center}
\begin{tabularx}{\textwidth}{@{}p{0.19\textwidth}X X@{}}
\toprule
Header group & Fields & Purpose \\
\midrule
format & magic, version, header length, value type, slot size &
reject unknown physical representations \\
shape & capacity, active length, layout hash, schema hash &
establish vector compatibility \\
publication & writer epoch, created and updated timestamps &
detect stable reads and identify data age \\
operations & writer flags and heartbeat &
distinguish liveness from value change \\
\bottomrule
\end{tabularx}
\end{center}

The current value types are dense \(f64\) and fixed quote, trade, and TAQ
records.  All shared records are fixed-size, pointer-free, naturally aligned,
and compatible with a C layout.  A schema hash identifies the payload meaning,
while a layout hash identifies the entity ordering.  Neither substitutes
for the other.

\section{The logical hypercube}

A collection of slices forms a logical cube even when no monolithic
multidimensional allocation exists.  For one generation,
\[
  X_g[u,f]=S_{f,g}[\,\iota(u)\,].
\]
The implemented engine names four axes directly: generation, entity, primitive
field, and calculated node.  Applications may add asset class, venue, factor,
family, horizon, universe, scenario, portfolio, or model version:
\[
  X[g,u,f,h,v,s,\ldots].
\]
Most of this space is not materialized.  A slice is dense along the entity
axis with the other coordinates fixed.  A catalog maps a semantic name such
as \texttt{signal} to the file, layout, type, and role that materialize one
such slice.

\begin{figure}[ht]
\centering
\begin{tikzpicture}[node distance=7mm and 8mm]
  \node[slicemeta] (layout) {\textbf{Layout}\\entity \(\leftrightarrow\) slot};
  \node[slicehot, right=of layout] (rows) {\textbf{Input rows}\\named primitive fields};
  \node[slicehot, right=of rows] (cube) {\textbf{Hypercube}\\field and linear nodes};
  \node[sliceaction, right=of cube] (outputs) {\textbf{Published slices}\\one per selected node};
  \draw[sliceflow] (layout) -- (rows);
  \draw[sliceflow] (rows) -- (cube);
  \draw[sliceflow] (cube) -- (outputs);
  \node[slicemeta, below=10mm of cube] (generation)
    {\textbf{Generation}\\mode and observation time};
  \draw[sliceflow] (generation) -- (cube);
  \draw[sliceflow] (generation) -- (outputs);
\end{tikzpicture}
\caption{Physical slices are aligned one-dimensional slices.  Layout and
generation metadata give the collection its hypercube meaning.}
\label{fig:cube-slices}
\end{figure}

This formulation avoids two extremes.  It does not flatten all meaning into an
untyped matrix, and it does not allocate every possible combination of
dimensions.  The catalog and calculation registry describe a sparse set of
materialized slices; consumers request only the axes and generations required
for their operation.

\section{Time, observation, and publication}

``Latest'' is not a sufficient temporal contract for financial decisions.  A
cell may carry:
\[
  \tau(c)=(t_{\mathrm{event}},t_{\mathrm{ingest}},
           t_{\mathrm{compute}},t_{\mathrm{publish}},g).
\]
Event time states when the source says the observation occurred.  Ingest time
states when the local system received it.  Compute and publish time bound the
derived value's latency.  Generation \(g\) states which coherent calculation
the value belongs to.

The current quote and trade records retain source and ingest timestamps.
TAQ records preserve the quote known locally when a trade was processed,
together with quote age and flags for missing, stale, or temporally inverted
context.  Hypercube \texttt{CellValue} records node, entity key, finite value,
and the oldest contributing observation time.  \texttt{NodeStatus} records
emitted and missing counts plus local compute time.  Snapshot time, slice
updated time, and writer heartbeat remain distinct.  An application may build
a freshness service level as a predicate, not an intuition:
\[
  \operatorname{fresh}(c,t_d)
  \iff \operatorname{status}(c)=\mathrm{ok}
  \land 0\le t_d-t_{\mathrm{publish}}(c)\le\Delta_{\max}
  \land t_{\mathrm{compute}}\le C_{\max}.
\]

\part{Consistency and Publication}

\section{What ``nonlocking read'' means}

The desired property is precise: a reader of an already mapped slice should
not acquire a process-shared mutex or wait for a writer-owned critical section.
This is a \emph{reader-lockless} contract, not automatically a formal
wait-free guarantee.  An optimistic reader may have to retry, and under a
continuous writer it may fail to obtain a stable version within its retry
budget.

Sequence counters are a standard form of this tradeoff.  The Linux
documentation describes lockless readers that accept data only when an even
sequence value is unchanged across the read; it also notes that readers can
spin while a writer is interrupted and that raw counters do not serialize
multiple writers \cite{seqlock}.  The current format consequently assumes one
writer per slice and bounds reader retries.

\section{Point-state protocol}

Quote, trade, and TAQ slices use one sequence counter per slot.  For slot
\(i\), the writer performs
\[
\begin{aligned}
  q_i&\leftarrow 2k+1 && \text{publication in progress},\\
  S[i]&\leftarrow v && \text{copy fixed payload},\\
  q_i&\leftarrow 2k+2 && \text{publication complete}.
\end{aligned}
\]
The reader loads \(q_0\), rejects an odd value, copies the payload, and loads
\(q_1\).  It accepts only if
\[
  q_0=q_1\quad\land\quad q_1\bmod 2=0.
\]
An update to AAPL therefore does not make a reader of MSFT retry.  This is the
appropriate granularity for independently changing latest-state records.

The point snapshot of an entire struct slice is a collection of individually
coherent cells.  It is not, by itself, proof that every slot represents the
same market instant.  Consumers that need an atomic universe snapshot require
a generation-level contract.

\section{Vector protocol}

Dense \(f64\) slices use one epoch for the vector.  A writer makes the epoch
odd, updates one or many cells, and publishes the next even epoch.  A reader
loads \(e_0\), scans or copies the vector, and loads \(e_1\); it accepts when
\[
  e_0=e_1\quad\land\quad e_1\bmod 2=0.
\]
Dot products validate both input epochs before and after the operation.
Layout, schema, and active-length checks precede the arithmetic.

This protocol gives a coherent view of one slice without a reader mutex.  Its
scope is limited:
\begin{itemize}
  \item a busy writer may force a reader to retry;
  \item separate slices have separate epochs, so sequentially stable copies may
        still belong to different logical input generations;
  \item one writer must own publication for a slice;
  \item memory-order behavior must be validated on every supported
        architecture and foreign-language reader implementation.
\end{itemize}

\section{Consistency classes}

Different data uses should pay for different consistency:
\begin{center}
\begin{tabularx}{\textwidth}{@{}p{0.15\textwidth}p{0.22\textwidth}X X@{}}
\toprule
Class & Example & Read contract & Failure behavior \\
\midrule
cell-latest & quote display, health probe &
one slot, per-cell sequence & retry or report unavailable \\
vector-stable & exposure, rank, dot product &
one slice, stable even epoch & bounded retry \\
generation-coherent & ETF value or another multi-field decision &
named immutable multi-slice generation & reject mixed generation \\
decision-grade & order or rebalance input &
coherent generation plus freshness, authority, and risk checks &
fail closed; no action \\
\bottomrule
\end{tabularx}
\end{center}

The table prevents a common design error: imposing a global lock on all state
because a small subset of decisions require a global snapshot.  Latest quote
display can tolerate independently published slots.  A cross-sectional rank
needs one stable vector.  An ETF fair-value calculation may need constituent
prices, currencies, holdings, and shares outstanding from one declared
generation.  A trade requires that coherence plus policy and authority.

\section{Published immutable generations}

The next publication layer is a small multi-version protocol.  A writer builds
generation \(g+1\) in inactive buffers while readers continue to use \(g\).
After every slice passes shape, schema, and completeness checks, the writer
publishes one compact manifest:
\[
  M_{g+1}=(g+1,t,\lambda,\{(f,e_f,h_f)\},H),
\]
where \(e_f\) and \(h_f\) are the final epoch and content identity of each
slice and \(H\) authenticates the manifest.  Publication changes a single
active-generation reference.  A reader pins \(M_g\), reads only the referenced
immutable buffers, and verifies that the active generation did not change
before accepting the result.

This is a read-copy-publish design related to multiversion and RCU techniques.
RCU separates removal from reclamation so readers can traverse an existing
version without acquiring a conventional lock \cite{rcu}; snapshot isolation
likewise gives a transaction a named database version \cite{snapshot}.  The
slice implementation still needs its own lifecycle rules:
\begin{itemize}
  \item a generation is immutable after publication;
  \item buffer reuse waits until no reader can reference the old generation;
  \item a slow reader is bounded or copied out so it cannot retain generations
        indefinitely;
  \item publication uses atomic rename or an aligned atomic generation word;
  \item crash recovery selects the last complete manifest, never half-written
        vectors.
\end{itemize}

No immutable multi-slice manifest or reader-pinning protocol is implemented in
version 1.  Consumers that need a coherent decision across several nodes
should use the in-process \texttt{Snapshot}; general multi-slice publication
remains an explicit extension and must not be inferred from independent slice
epochs.

\begin{figure}[ht]
\centering
\begin{tikzpicture}[node distance=8mm and 10mm]
  \node[slicehot] (active) {\textbf{Generation \(g\)}\\immutable, reader-visible};
  \node[slicemeta, right=of active] (build) {\textbf{Build \(g+1\)}\\inactive buffers};
  \node[slicemeta, right=of build] (validate) {\textbf{Validate}\\shape, schema, coverage};
  \node[sliceaction, right=of validate] (publish) {\textbf{Publish}\\one generation marker};
  \draw[sliceflow] (active) -- node[above,font=\scriptsize\sffamily]{read} (build);
  \draw[sliceflow] (build) -- (validate);
  \draw[sliceflow] (validate) -- (publish);
  \draw[slicefuture] (publish.south) to[bend right=27]
    node[below,font=\scriptsize\sffamily]{new readers pin \(g+1\)}
    (active.south);
\end{tikzpicture}
\caption{Copy-publish generations keep readers on immutable state while the
next coherent set of slices is assembled.}
\label{fig:generation}
\end{figure}

\section{Sharding and progress}

Double buffering every vector is attractive when the working set is small.
For larger or faster-changing cubes, the publication unit can be a shard:
entity range, field group, asset class, or semantic family.  A generation
manifest then references a vector of immutable shard versions.  Copy-on-write
touches only changed shards while the manifest preserves a coherent global
view.

Progress must be measured rather than assumed.  Relevant counters include
writer duration, epoch changes per second, reader retries, snapshot failures,
generation publish latency, oldest pinned generation, and reclaimed bytes.
If retry tails grow under ingestion load, the answer is not an unbounded spin
loop.  The system can use smaller shards, longer immutable generations, a
copied snapshot, or a stale-but-declared generation according to the
consumer's consistency class.

\part{Hypercube Execution}

\section{Runtime contract}

The current Rust engine receives
\[
  U_g=(g,t,m,R,N),
\]
where \(g\) is a strictly increasing generation, \(t\) is its Unix-millisecond
observation time, \(m\) is the execution mode, \(R\) contains keyed
\texttt{InputRow} values and named primitive fields, and \(N\) is the complete
\texttt{NodeSpec} graph.  It emits
\[
  H_g=(g,t,m,n,V,Q),
\]
where \(n\) is the entity count, \(V\) contains ordered
\texttt{CellValue} records, and \(Q\) contains ordered \texttt{NodeStatus}
diagnostics.

The modes are \texttt{Live}, \texttt{Replay}, and \texttt{Batch}.  They are
recorded context, not separate calculation implementations.  A replay changes
the supplied clock and rows rather than the meaning of a node.  The engine
rejects a generation that does not advance beyond the last successful update.

\section{Calculation nodes}

The implemented engine has two node kinds and five cross-sectional transforms:
\begin{center}
\begin{tabularx}{\textwidth}{@{}p{0.22\textwidth}X X@{}}
\toprule
Node & Input and operation & Output \\
\midrule
\texttt{Field} & one named primitive field from every entity row &
aligned raw values before the selected transform \\
\texttt{Linear} & required or optional upstream nodes and declared weights &
weighted value per entity, optionally normalized by available absolute weight \\
\midrule
\multicolumn{3}{@{}l@{}}{\texttt{Identity}, \texttt{ZScore},
\texttt{Rank}, \texttt{Percentile}, and \texttt{RankZScore}} \\
\bottomrule
\end{tabularx}
\end{center}

Node declaration order is not execution order.  The engine resolves
dependencies topologically and returns values in declaration order.  Duplicate
rows or nodes, empty identifiers, duplicate inputs, non-finite weights, cycles,
and unknown required dependencies fail the update.  A missing required value
suppresses that entity's linear output; an optional dependency may be absent.

\section{Cross-sectional algebra}

For node \(j\) and entity \(i\), let \(x_{ij}\) be a raw current value.
The engine can publish raw values, ordinary z-scores, ranks, percentiles, or
rank z-scores.  A simple linear composite is
\[
  c_i = \mathcal N\!\left(
    \frac{\sum_{j\in A_i}w_jx_{ij}}
         {\sum_{j\in A_i}|w_j|}
  \right),
\]
where \(A_i\) is the set of available dependencies for entity \(i\) and
\(\mathcal N\) is the declared output transform.  The denominator is used only
when \texttt{normalize\_weights} is true; otherwise the raw weighted sum is
transformed.  A zero available denominator emits no value.

\section{Resident state and incremental work}

The engine retains only the last successfully completed generation as a
monotonicity guard.  Every accepted update supplies complete rows and node
specifications and recalculates the graph; there are no built-in price
histories, session accumulators, alias registries, or cross-generation
incremental cache.

The natural extension is change-directed evaluation:
\[
  \Delta X_g \longrightarrow
  \{\phi:\operatorname{deps}(\phi)\cap\Delta X_g\ne\varnothing\}
  \longrightarrow \Delta Y_g.
\]
In a future incremental engine, only nodes whose dependency cones intersect
changed cells would need to run.
Logical generation and event time must remain part of the cache key; otherwise
a fast cache can quietly return a value computed from the wrong observation
boundary.

\section{Hot state, persistence, and transport}

\begin{figure}[ht]
\centering
\begin{tikzpicture}[node distance=7mm and 7mm]
  \node[slicemeta] (input) {\textbf{Caller or injector}\\complete update};
  \node[slicehot, right=of input] (cube) {\textbf{HypercubeEngine}\\validate and evaluate};
  \node[slicehot, right=of cube] (snapshot) {\textbf{Snapshot}\\values and status};
  \node[sliceaction, right=of snapshot] (publish) {\textbf{Outputs}\\slices, HTTP/SSE, or terminal};
  \draw[sliceflow] (input) -- (cube);
  \draw[sliceflow] (cube) -- (snapshot);
  \draw[sliceflow] (snapshot) -- (publish);
\end{tikzpicture}
\caption{The implemented public path.  Vendor feeds, durable persistence,
replication, and production gateways are outside the repository.}
\label{fig:system}
\end{figure}

\texttt{SlicePublisher} writes one configured \(f64\) slice per selected node
and publishes \texttt{layout.json} and \texttt{catalog.json}.  Those files are
independently coherent; the in-process \texttt{Snapshot} is the coherent
whole-generation object.  The synthetic server also serves a dependency-free
browser dashboard, a JSON snapshot, and an SSE stream.  The \texttt{etf} and
\texttt{pairs} examples run separate top-like terminal monitors over basket
premiums and cointegrating residuals.  These paths make the architecture
inspectable without external services.

The version-1 file ABI and Rust readers are a public local API, not a remote
service contract.  The example server has no authentication, entitlement,
backpressure, or production reconnect guarantees.  Vendor feeds, durable
stores, host-to-host transport, and a governed gateway require separate
adapters and policies; none is hidden in the public engine.

\part{Composable Cells}

\section{From values to functions}

The current cube materializes primitive values and runs a fixed set of node
types.  The general cell model is
\[
  c=(\tau,\sigma,\pi,\nu),
\]
where \(\tau\) is a value type, \(\sigma\) is shape and dimension metadata,
\(\pi\) is provenance and publication state, and \(\nu\) is either a
materialized value or a typed expression.

For a function cell,
\[
  c_j := \phi_j(c_{i_1},\ldots,c_{i_k};\theta_j),
\]
the registry stores the function identity and version, typed inputs,
parameters \(\theta_j\), output shape, missing-data policy, and permitted
state.  The dependency graph is explicit.  Acyclic pure expressions can be
topologically evaluated; windowed or recursive expressions require declared
state and logical-time semantics.

Composability comes from retaining the same cell contract at each stage.  A
quote midpoint can feed a return; returns can feed factors; factors can feed a
portfolio or index; the resulting discrepancy can feed a decision rule.  Each
output remains addressable, versioned, and independently inspectable.

\section{Function classes}

Not every function should have the same authority:
\begin{center}
\begin{tabularx}{\textwidth}{@{}p{0.17\textwidth}p{0.25\textwidth}X@{}}
\toprule
Class & Examples & Contract \\
\midrule
pure scalar/vector & midpoint, return, dot product, normalization &
deterministic for a generation; no hidden I/O \\
stateful analytic & VWAP, rolling covariance, session range &
declared state, clock, retention, reset, and replay semantics \\
model inference & learned score, volatility estimate &
versioned artifact, features, calibration, and resource limits \\
decision & threshold, rebalance proposal, hedge intent &
produces an intent and explanation; no execution authority \\
effect adapter & send order, alert, persist, publish &
explicit capability, idempotency, audit, timeout, and failure policy \\
\bottomrule
\end{tabularx}
\end{center}

The distinction between a decision and an effect is structural.  A function
cell may calculate that a discrepancy exceeds a threshold.  It should not
silently turn a read of that cell into an order.  Effect adapters consume
published decision intents only after checking generation, freshness,
entitlement, position limits, market state, and a unique action key.

\section{ETF and index example}

Consider an ETF \(e\) with constituent quantities \(q_{ei}\), local prices
\(p_i(t)\), currency conversion rates \(x_i(t)\), cash and accrued value
\(C_e(t)\), liabilities \(L_e(t)\), and shares outstanding \(N_e(t)\).  A
simplified computed value per share is
\[
  \widehat V_e(t)=
  \frac{C_e(t)-L_e(t)+\sum_i q_{ei}p_i(t)x_i(t)}
       {N_e(t)}.
\]
An index uses analogous constituent arithmetic with its declared weights,
corporate-action rules, and divisor.  The observed ETF midpoint \(m_e(t)\)
gives a discrepancy
\[
  \delta_e(t)=m_e(t)-\widehat V_e(t),\qquad
  \delta^{\mathrm{bps}}_e(t)=
  10^4\frac{m_e(t)-\widehat V_e(t)}{\widehat V_e(t)}.
\]

In the cube, each input is an aligned slice or a declared keyed join:
\begin{align*}
  P_g &= \text{constituent price vector},\\
  X_g &= \text{currency conversion vector},\\
  Q_{e,g} &= \text{holdings vector},\\
  \widehat V_{e,g} &=
    \phi_{\mathrm{nav}}(P_g,X_g,Q_{e,g},C_{e,g},L_{e,g},N_{e,g}),\\
  \delta_{e,g} &= \phi_{\mathrm{delta}}(m_{e,g},\widehat V_{e,g}).
\end{align*}
The calculation is accepted only when price, FX, holdings, cash, liabilities,
and share count satisfy their observation and generation policies.  Holdings
from last night's file may be valid with an explicit as-of time; a stale FX
rate may not be.  The cell graph makes those choices visible.

\subsection{Runnable ETF monitor}

The checked-in \texttt{etf} example implements a compact version of the basket
calculation.  It creates stable synthetic constituents and long-only baskets
whose weights sum to one.  For ETF \(j\), it forms a basket return and advances
fair value along the constituent axis:
\[
  r^{\mathrm{NAV}}_{j,g}=\sum_{i=1}^{N} w_{ji}r_{i,g},
  \qquad
  V_{j,g}=V_{j,g-1}(1+r^{\mathrm{NAV}}_{j,g}).
\]
The simulated ETF price is \(P_{j,g}=V_{j,g}(1+u_{j,g})\), where the fractional
premium follows
\[
  u_{j,g}=\phi u_{j,g-1}+\sigma\eta_{j,g},
  \qquad \phi=0.88,\quad \sigma=0.00040.
\]
The terminal reports premium in basis points and against the stationary model
standard deviation:
\[
  q_{j,g}=10^4\frac{P_{j,g}-V_{j,g}}{V_{j,g}},
  \qquad
  z_{j,g}=\frac{u_{j,g}}{\sigma/\sqrt{1-\phi^2}}.
\]
Hypercube publishes ETF-aligned fair values, prices, premiums, model
\(z\)-scores, and cross-sectional premium ranks.  This keeps the two axes
explicit: basket arithmetic reduces a constituent vector to one ETF value; a
Hypercube node produces another aligned ETF vector.

\subsection{Runnable pairs monitor}

The \texttt{pairs} example gives each pair a stable identity and two price
paths.  Leg \(B\) carries a stochastic trend, while leg \(A\) shares it through
hedge ratio \(\beta_j\).  Their cointegrating residual is
\[
  y_{j,g}=\log A_{j,g}-\alpha_j-\beta_j\log B_{j,g},
  \qquad
  y_{j,g}=\phi_j y_{j,g-1}+\sigma_j\eta_{j,g}.
\]
The graph publishes current prices, the standardized residual, model
half-life, and a cross-sectional rank of absolute residual size:
\[
  z_{j,g}=\frac{y_{j,g}}{\sigma_j/\sqrt{1-\phi_j^2}},
  \qquad
  h_j=\frac{\log(1/2)}{\log\phi_j}.
\]
The table shows the largest live dislocations and the sign-implied hedged
position.  Fixed parameters keep the example readable; estimation,
stationarity tests, structural breaks, costs, borrow, and exposure controls
remain outside it.

\begin{figure}[ht]
\centering
\resizebox{\textwidth}{!}{%
\begin{tikzpicture}[node distance=6mm and 8mm]
  \node[slicehot] (px) {constituent prices};
  \node[slicehot, below=of px] (fx) {FX conversion};
  \node[slicemeta, below=of fx] (hold) {holdings / divisor};
  \node[slicehot, right=16mm of fx] (fair) {computed ETF/index value};
  \node[slicehot, above=of fair] (obs) {observed market value};
  \node[sliceaction, right=of fair] (delta) {delta and diagnostics};
  \node[sliceaction, right=of delta] (intent) {bounded trade intent};
  \node[slicemeta, right=of intent] (effect) {authorized executor};
  \draw[sliceflow] (px.east) -- (fair.west);
  \draw[sliceflow] (fx.east) -- (fair.west);
  \draw[sliceflow] (hold.east) -- (fair.west);
  \draw[sliceflow] (fair) -- (delta);
  \draw[sliceflow] (obs.east) -- (delta.west);
  \draw[slicefuture] (delta) -- (intent);
  \draw[slicefuture] (intent) -- (effect);
\end{tikzpicture}%
}
\caption{The calculation is already vector-shaped.  Dashed edges mark future
decision and effect layers, not current slice behavior.}
\label{fig:etf}
\end{figure}

The same pattern supports baskets, hedges, factor-mimicking portfolios, risk
decompositions, and observed-versus-model residuals.  The significant
abstraction is not the ETF formula itself.  It is the ability to compose
versioned values and functions while preserving a coherent generation and an
explicit effect boundary.

\section{Action boundary}

A future trading adapter should accept an immutable intent:
\[
  I=(g,\text{strategy},\text{instrument},q,\text{limit},
     \text{expiry},\text{evidence},k),
\]
where \(k\) is an idempotency key.  Before submission it must verify:
\begin{enumerate}
  \item generation \(g\) is decision-grade and within its freshness budget;
  \item the calculation, model, and policy versions are approved;
  \item the account has authority for the instrument and action;
  \item position, notional, liquidity, price, and loss limits pass;
  \item the market is in an allowed session and no kill switch is active;
  \item action key \(k\) has not already been consumed.
\end{enumerate}

The executor records acknowledgement, rejection, fill, and cancellation events
back into the evidence path.  These events may become new cube inputs, but the
dependency must not create an unbounded accidental feedback loop.  Stateful
control policies need explicit cycle semantics, rate limits, and termination
conditions.

\part{Implementation and Evidence}

\section{Current implementation}

The repository contains two principal Rust crates:
\begin{center}
\begin{tabularx}{\textwidth}{@{}p{0.22\textwidth}X X@{}}
\toprule
Component & Implemented responsibilities & Present boundary \\
\midrule
\texttt{hypercube-slice} &
layout and catalog validation; versioned headers; \(f64\), quote, trade, and
TAQ mappings; point and vector readers/writers; stable snapshots, sums, dot
products, and top-absolute scans &
little-endian; one writer per slice; local file-backed mappings \\
\texttt{hypercube-engine} &
generation-checked updates; field and linear nodes; five cross-sectional
transforms; deterministic snapshots and status; output
\texttt{SlicePublisher}; generic and OU market injectors; HTTP/SSE plus
\texttt{etf} and \texttt{pairs} terminal demonstrations &
complete in-process rows per update; no vendor adapter, persistence, general
function cells, or production web service \\
\bottomrule
\end{tabularx}
\end{center}

The slice crate supplies fixed quote, trade, and TAQ record schemas but does
not ingest a live feed.  Callers own the meanings and provenance of records
they publish.  The engine accepts ordinary in-process rows rather than reading
input slices directly.  Its publisher projects selected calculated nodes into
\(f64\) slices sharing one entity layout; it does not publish factor-family
rollups, portfolio projections, or application-specific summaries.

\section{Scale model and reproducible baseline}

For \(F\) published nodes and \(N\) entities, dense \(f64\) payloads use
approximately
\[
  8FN\ \text{bytes}
\]
plus one 256-byte header per slice.  The checked-in semi-live smoke run uses
eight entities and five nodes.  It emits 40 cell values and five 320-byte
slices:
\[
  256 + 8\cdot 8 = 320\ \text{bytes per slice}.
\]
The dashboard, JSON snapshot, and SSE stream expose that same state.  This is
a functional and packaging baseline, not a latency or throughput benchmark;
the public repository makes no development-host performance claim.

\section{Conformance tests}

The current workspace has thirteen Rust tests:
\begin{enumerate}
  \item six engine, publisher, and synthetic-injector tests cover dependency
        cycles, stale generations, tie-aware ranking, deterministic generic and
        OU market input, and aligned publication;
  \item seven slice tests cover layouts and catalogs, format compatibility,
        \(f64\) vector reads and dot products, heartbeat semantics, and quote,
        trade, and TAQ records.
\end{enumerate}

Foreign-language byte fixtures, high-contention stress, interrupted writers,
cross-process x86-64 and aarch64 evidence, immutable generation pinning, and
crash recovery remain conformance work.  Performance reports should state
entity count, slice count, update
distribution, writer rate, reader count, operation, page residency, CPU
architecture, flush policy, and percentile methodology.  Useful
measurements include point-read latency, stable-vector latency, retries per
accepted read, generation publish time, changed-shard bytes, and end-to-end
input-to-node-to-publication latency.

\section{Limits}

The current implementation establishes a credible vertical slice, not a
finished universal memory system:
\begin{itemize}
  \item published output vectors are stable per slice but not globally atomic
        across the complete output set;
  \item vector readers use optimistic retry and can fail to progress within a
        fixed budget under sustained writes;
  \item shared atomic and payload behavior requires explicit cross-process,
        cross-language, x86-64, and aarch64 conformance evidence;
  \item layouts are stable while running; dynamic mutation requires a new
        layout and coordinated publication;
  \item replication and general live gateway generation protocols remain
        separate work;
  \item vendor ingestion, durable persistence, and direct input-slice adapters
        are outside the public repository;
  \item the hypercube is a typed calculation runtime, but arbitrary function
        cells, automatic incremental invalidation, model sandboxing, and GPU
        lowerings are not implemented;
  \item action cells and trading authority are future architecture and must
        not be implied by the presence of a calculated discrepancy.
\end{itemize}

These limits are part of the contribution.  They identify where a simple
memory-mapped vector is sufficient and where stronger publication, lifecycle,
or authority semantics are required.

\section{Open-source research sequence}

Starting from the current public version, development can proceed in bounded
stages:
\begin{enumerate}
  \item stabilize the published format, layout and catalog schemas, Rust
        implementation, synthetic fixture, and current conformance tests;
  \item add portable readers and an ABI corpus that proves the same bytes have
        the same meaning outside Rust;
  \item benchmark point counters and vector epochs under declared writer and
        reader loads on x86-64 and aarch64;
  \item implement a general immutable generation manifest with reader pinning,
        crash recovery, and optional changed-shard publication;
  \item formalize typed pure and stateful function-cell interfaces, dependency
        identity, change propagation, and replay equivalence;
  \item build the ETF/index example through fair value and discrepancy as a
        deterministic reference application;
  \item specify decision intents and a simulated effect adapter before any
        connection to live execution.
\end{enumerate}

Each stage should be useful independently.  A community member may only need a
fast stable-layout shared vector, a factor researcher may use the replay/live
hypercube contract, and a systems researcher may focus on multi-slice
generation consistency.  The project does not need to expose trading authority
to make its memory and composition model valuable.

\section{Conclusion}

\SLICE{} reduces current analytical state to a small, inspectable
contract: fixed typed cells, stable entity slots, versioned headers, and
reader-lockless publication protocols.  The hypercube restores the dimensions
that a raw vector alone cannot express---entity, primitive field, calculated
node, and generation---while keeping common operations compatible with
ordinary linear algebra.  Applications may add factor, horizon, universe, and
provenance coordinates without changing the version-1 core.

The concurrency design is intentionally plural.  Per-slot counters serve
independently changing point state and slice epochs serve stable vector scans.
Immutable generation manifests are the proposed stronger class for coherent
multi-vector decisions.  This
avoids locking an entire analytical universe for every market change without
pretending that all consumers can accept the same consistency.

Function cells would extend the model from shared values to a richer shared
computation graph.  Pure and explicitly stateful nodes could compose data into
returns, factors, portfolios, ETF or index values, and observed-model
residuals.  Decision nodes may produce bounded intents, but effects remain
behind a separately governed adapter.  That separation lets composability grow
without allowing convenience to erase provenance, freshness, or authority.

The germinal system is therefore already present: a live memory plane, an
aligned multidimensional calculation model, and versioned outputs that can be
read repeatedly at low cost.  The open-source work is to make its guarantees
portable, its generation semantics complete, its benchmarks reproducible, and
its function boundary explicit.

\begin{thebibliography}{9}

\bibitem{posixmmap}
The Open Group,
\emph{POSIX.1-2024: mmap---Map Pages of Memory},
\url{https://pubs.opengroup.org/onlinepubs/9799919799/functions/mmap.html}.

\bibitem{arrow}
Apache Arrow Project,
\emph{Arrow Columnar Format},
\url{https://arrow.apache.org/docs/format/Columnar.html}.

\bibitem{scidb}
M. Stonebraker, P. Brown, D. Zhang, and J. Becla,
\emph{SciDB: A Database Management System for Applications with Complex
Analytics},
Computing in Science \& Engineering, 15(3), pp.~54--62, 2013,
\url{https://doi.org/10.1109/MCSE.2013.19}.

\bibitem{blas}
C. L. Lawson, R. J. Hanson, D. R. Kincaid, and F. T. Krogh,
\emph{Basic Linear Algebra Subprograms for FORTRAN Usage},
ACM Transactions on Mathematical Software, 5(3), pp.~308--323, 1979,
\url{https://doi.org/10.1145/355841.355847}.

\bibitem{seqlock}
Linux Kernel Documentation,
\emph{Sequence Counters and Sequential Locks},
\url{https://www.kernel.org/doc/html/latest/locking/seqlock.html}.

\bibitem{rcu}
Linux Kernel Documentation,
\emph{RCU Concepts},
\url{https://docs.kernel.org/RCU/rcu.html}.

\bibitem{snapshot}
H. Berenson, P. Bernstein, J. Gray, J. Melton, E. O'Neil, and P. O'Neil,
\emph{A Critique of ANSI SQL Isolation Levels},
SIGMOD, pp.~1--10, 1995,
\url{https://doi.org/10.1145/223784.223785}.

\bibitem{adaptive}
U. A. Acar, G. E. Blelloch, and R. Harper,
\emph{Adaptive Functional Programming},
ACM Transactions on Programming Languages and Systems, 28(6),
pp.~990--1034, 2006,
\url{https://doi.org/10.1145/1186632.1186634}.

\bibitem{differential}
F. McSherry, D. G. Murray, R. Isaacs, and M. Isard,
\emph{Differential Dataflow},
CIDR, 2013,
\url{https://www.cidrdb.org/cidr2013/Papers/CIDR13_Paper111.pdf}.

\end{thebibliography}

\end{document}
